\chapter{Problem Analysis}
\section{ROS-Based Robotic Systems Overview}
Nowadays robotic systems are increasingly built as distributed software architectures composed of multiple interacting components.
The Robot Operating System (ROS) has become the standard framework when developing such systems, providing a flexible middleware
that enables modularity, reuse, and interoperability across heterogeneous hardware and software platforms. In ROS-based systems,
functionality is decomposed into nodes, each responsible of a well-defined task. These nodes communicate through message passing mechanisms using topics, services, and actions, forming a 
dynamic computation graph. This publisher-subscriber architecture enables developers to construct complex robotic behaviors. Nodes can be developed independently, deployed on different machines, and connected at runtime,
allowing robotic systems to scale from single embedded devices to distributed setups.

Additionally, ROS supports a wide ecosystem of third-party packages that implement common algorithms and hardware interfaces,
significantly reducing development effort. However, the flexibility offered by ROS also introduces complexity. The overall system
behavior emerges not only from the implementation of individual nodes, but also from how they are connected, configured, and parameterized.
The ROS computation graph may involve dozens of nodes with intricate dependencies on sensors, coordinate frames, timing assumptions, and
runtime parameters. As a result, the correctness and performance of the system depend heavily on the configuration choices made during deployment.
More than that, ROS itself intentionally avoids imposing strong constraints on how systems should be structured. While this design choice promotes
generality and reusability, it places a significant burden on developers to ensure that node interactions are valid for a given task.

As robotic systems grow in scale and complexity, understanding and managing these interactions becomes challenging,
particularly when systems must be adapted to new missions, hardware configurations, or operational environments. 
Beyond software modularity, ROS-based systems must also operate under real-world constraints that further complicate system design. Robotic
applications are often required to integrate heterogeneous sensors, operate under strict timing requirements, and interact with uncertain
and dynamic environments. These constraints introduce implicit assumptions into system configurations, such as expected sensor frequencies,
coordinate frame hierarchies, and synchronization between perception and control pipelines. While ROS provides the communication primitives
needed to express these interactions, it does not enforce correctness or consistency at the system level. As a result, many design decisions 
are encoded implicitly in configuration files rather than being explicitly represented or validated. This implicit knowledge becomes increasingly
difficult to maintain as systems evolve, particularly when configurations are reused across different platforms or mission scenarios.

\section{ROS Configuration and Orchestration Challenges}
Configuring a ROS-based robotic system involves more than launching executable nodes. Developers must specify a wide range of parameters that govern
node behavior, communication interfaces, sensor usage, and algorithmic settings. These configurations are typically expressed through a combination 
of launch files, parameter files, often in YAML format, and environment-specific settings. Even minor changes, such as replacing a sensor or adjusting
an algorithm parameter, may require coordinated updates across multiple configuration files. Orchestration further complicates this process. In 
distributed deployments, nodes may be executed across multiple machines or containers, requiring careful management of networking, synchronization, 
and lifecycle events. Ensuring that nodes are launched in the correct order, recover gracefully from failures, and remain consistent across restarts 
is non-trivial, particularly when systems are deployed in dynamic or resource-constrained environments. 

Existing orchestration mechanisms within ROS provide only limited support for these challenges. Launch files primarily describe static execution
graphs and offer minimal capabilities for conditional logic or dynamic adaptation. External orchestration frameworks, such as container-based 
deployment platforms, can manage process lifecycles and resource allocation but operate largely independently of ROS-specific semantics and in most 
cases such system are not designed and do not support ROS components. As a result, developers must manually bridge the gap between system-level
orchestration and ROS-level configuration. Another key challenge lies in the lack of abstraction over mission intent. ROS configurations describe 
how components are connected and parameterized, but not why a particular configuration is appropriate for a given task. Mission objectives, such as 
exploration, mapping, or navigation are implicitly encoded in configuration choices rather than explicitly represented. This makes it difficult to
reason about configurations at a higher level or to systematically adapt systems to new tasks without extensive manual intervention.

\section{Limitations of Manual and Static Approaches}
The predominant approach to configuring ROS-based systems remains manual and largely static. Developers typically create configuration files by hand,
relying on experience, documentation, and trial-and-error to achieve correct system behavior. While this approach may be manageable for small systems 
or fixed use cases, it does not scale well as systems grow in complexity or must support multiple operational modes. Manual configuration is inherently
error-prone. Parameter mismatches, incompatible topic connections, or incorrect assumptions about sensor availability can lead to subtle runtime 
failures that are difficult to diagnose. Moreover, configurations are often tightly coupled to specific hardware setups and mission contexts, reducing
their reusability across different scenarios. Adapting an existing system to a new task frequently involves duplicating and modifying large portions 
of configuration code, increasing maintenance effort and the risk of inconsistencies. Static configurations also limit adaptability. Once a system is
launched, its structure and parameters typically remain fixed, even if environmental conditions or task requirements change. While dynamic 
reconfiguration mechanisms exist, they are rarely used systematically due to their complexity and lack of integration with higher-level reasoning.
As a result, robotic systems often lack the ability to adjust their behavior automatically in response to changing mission objectives. These limitations
show a fundamental gap between the low-level configuration mechanisms provided by ROS and the high-level reasoning required for flexible, 
mission-driven robotic systems. Closing this gap requires methods that can reason about system structure, parameters, and dependencies in a 
more abstract manner, while still producing concrete, executable configurations.

The lack of standardized abstraction mechanisms further increases orchestration challenges in ROS-based systems. Configuration files often mix 
concerns related to hardware setup, algorithm selection, and mission-specific behavior, resulting in tightly coupled specifications that are 
difficult to reason about independently. While parameter servers and launch-time arguments provide some degree of flexibility, they do not offer
a systematic way to capture higher-level relationships between components. Consequently, orchestration logic is frequently implemented through ad-hoc
scripts or duplicated configuration fragments, making systems harder to extend and validate. This fragmentation also affects reproducibility, as small,
undocumented changes in configuration can significantly alter system behavior without being reflected in versioned source code. 

Another flaw is that understanding the impact of a single parameter change often requires detailed knowledge of multiple interacting components, as well as their 
underlying assumptions. This makes configuration tasks difficult to delegate or automate and limits the ability of non-expert users to deploy or 
modify robotic systems safely. Moreover, the absence of explicit validation or reasoning mechanisms means that many configuration errors are only 
discovered at runtime, potentially leading to system downtime or unsafe behavior. These factors underscore the need for approaches that can externalize
configuration knowledge, reduce manual effort, and support systematic reasoning about system structure and behavior.